\documentclass[11pt]{article}
\usepackage[review]{acl}
\usepackage{times}
\usepackage{latexsym}
\usepackage[T1]{fontenc}
\usepackage[utf8]{inputenc}
\usepackage{microtype}
\usepackage{inconsolata}
\usepackage{booktabs}
\usepackage{url}

\setlength{\emergencystretch}{2em}
\newcommand{\code}[1]{\path{#1}}
\newcommand{\metricbar}[1]{\rule{#1em}{0.8ex}}

\title{Explanatory Graph Context Rehydration for Agentic Systems:\\Bounded Retrieval Across Pull and Event-Driven Runtimes}

\author{Anonymous ACL submission}

\begin{document}
\maketitle

\begin{abstract}
Agent systems often assemble context through runtime-specific prompts,
application endpoints, and ad hoc caches, making context management difficult
to reuse, audit, or evaluate independently of execution. We present
\texttt{rehydration-kernel}, a graph-native context service that reconstructs
bounded operational context from projected state and exposes it through query
and asynchronous interfaces. The kernel models context as nodes with optional
detail plus typed relationship explanations that preserve why a downstream node
exists, including rationale, motivation, method, and decision linkage. We
evaluate the system with container-backed end-to-end workflows and four
explanatory use cases: failure diagnosis with rehydration-point recovery,
reconstruction of why a task was implemented in a particular way, interrupted
handoff with resumable execution, and constraint-preserving retrieval under
token pressure. Across these use cases, full explanatory context achieves
\mbox{1.0} explanation roundtrip fidelity and \mbox{1.0} causal
reconstruction in the full-detail setting, while retaining \mbox{1.0} causal
reconstruction under a \mbox{192}-token budget in the constraint-preserving
case. Removing node detail preserves explanation fidelity but reduces causal
reconstruction to \mbox{0.857}, while replacing explanatory relations with
structural edges reduces causal reconstruction to \mbox{0.143}, eliminates
rehydration-point recovery, and fails to preserve the dominant reason under
budget pressure. These results suggest that explanatory relationships carry
the dominant signal for restartable and auditable agent workflows, while node
detail primarily improves completeness. An LLM-in-the-loop evaluation with
GPT-5.4 and Claude Opus~4 across 18 scale--domain--relation configurations
confirms that explanatory context achieves 94\% on task, restart, and rationale
metrics versus 61\% for structural-only bundles, with the gap widening at
larger graph scales.
\end{abstract}

\section{Introduction}

Context management remains an under-specified systems problem in agent
architectures. In many practical stacks, the runtime that plans or executes
actions also assembles its own context from prompts, application APIs, caches,
and task-local heuristics. This coupling makes context assembly difficult to
reuse across runtimes, difficult to audit after the fact, and difficult to
evaluate as a system component in its own right.

This paper explores a different separation of concerns. With
\texttt{rehydration-kernel}, context rehydration is isolated into a reusable
kernel that reconstructs bounded graph context from projection state. Runtimes
consume already rehydrated context through stable query or event boundaries
instead of rebuilding that context inside each execution loop.

Our central claim is not merely that graph structure can hold context. It is
that explanatory relationships materially improve the usefulness of rehydrated
context for agent workflows. In our model, a relationship does not only encode
connectivity between two nodes. It may also preserve why the downstream node
exists, what decision produced it, what motivation drove it, or what method
verified it.

This distinction matters for operational questions that arise in realistic
agent systems: why was a task implemented in a particular way, which
relationships led to an incorrect implementation, and from which upstream node
should the system rehydrate and retry.

The paper makes three contributions. First, it presents a reusable kernel
architecture for bounded graph context rehydration over projected state.
Second, it introduces an explanatory relationship model that preserves
transition-level intent without moving domain ownership into the kernel.
Third, it provides a reproducible artifact showing that explanatory relations
contribute more to causal reconstruction and recovery than node detail alone.

\section{Problem Statement}

An agent needs bounded, actionable context, not raw graph storage. A raw graph
alone does not define what part of the graph should be selected, what
explanatory fields should survive retrieval, how to preserve detail under
token pressure, or how to present restartable reasoning context to the
runtime.

We therefore define the problem as follows: given projected graph state and
optional node detail, construct a bounded context bundle that preserves enough
local structure and explanation for an agent runtime to act, diagnose,
justify, or recover.

The design constraints in this repository are deliberate. The kernel must
remain generic and avoid product-specific nouns, support both pull and
event-driven entrypoints, avoid inventing semantic motivation on its own, and
preserve explanations emitted by an agent or domain application.

\section{Context Model}

The context model has four primitives: a root node that anchors retrieval,
neighbor nodes that define the selected local graph, relationships between
nodes, and optional extended node detail stored separately from the graph.

The model becomes stronger when relationships carry typed explanation. In the
current implementation, representative explanation fields include
\code{semantic_class}, \code{rationale}, \code{motivation}, and
\code{method}. The model also preserves \code{decision_id},
\code{caused_by_node_id}, \code{evidence}, \code{confidence}, and
\code{sequence}. This is an application-supplied explanation of transition
between nodes. The kernel preserves it, stores it, rehydrates it, and renders
it. The kernel does not infer it.

That split matches the domain boundary we want. The domain that knows why a
task exists or why a decision was taken is not the kernel. The kernel is the
infrastructure that keeps that explanation intact across projection, storage,
query, and rendering boundaries.

\section{System Architecture}

The repository implements the model as a graph-native context service with
async projection inputs over NATS, graph state in Neo4j, extended node detail
in Valkey, and query serving over gRPC. Runtime execution remains outside the
kernel boundary.

This architecture enables the same rehydrated context to be consumed by
pull-driven runtimes that query context before execution and event-driven
runtimes triggered from generated context bundles. The current quantitative
paper harness focuses on explanatory use cases, while broader integration
suites in the repository demonstrate cross-runtime and event-driven
operability.

\section{Related Work}

Our work is closest to four lines of prior work, but differs from each in
system boundary and objective.

First, GraphRAG-style systems use graph structure to organize retrieval for
language models. Edge et al. cast GraphRAG as graph-based query-focused
summarization over large corpora~\citep{edge2024localglobal}. That work is
close in its use of graph-shaped context, but it is document-centric. Our
kernel instead rehydrates projected operational state and serves it as runtime
context rather than constructing a graph from retrieved text.

Second, event-graph and narrative-graph approaches preserve richer temporal and
causal structure than plain knowledge graph triples. Our explanatory
relationship model follows that direction by preserving why a downstream node
exists rather than only which node comes next. ExCAR is especially relevant
because it demonstrates the value of event-graph evidence for explainable
causal reasoning~\citep{du2021excar}.

Third, hyper-relational knowledge graph research shows that relation qualifiers
carry contextual signal that plain triples lose. HAHE and HOLMES reinforce
this point for representation learning and multi-hop reasoning over qualified
relations~\citep{luo2023hahe,panda2024holmes}. Text2NKG extends the same
direction by supporting n-ary and event-based schemas that more closely match
how operational workflows are described in practice~\citep{luo2024text2nkg}.
Our contribution is to apply this principle in a systems setting where
qualified relationships must survive projection, storage, query, and rendered
context boundaries.

Finally, there is a semantic-modeling precedent for separating contextual and
conceptual information. GKR uses a layered graph representation with a
dedicated contextual layer rather than collapsing all semantics into a single
flat structure~\citep{kalouli2018gkr}. Our model is not a semantic parser, but
the split is aligned: node identity and node detail remain distinct from
relationship-level explanation about why a transition occurred.

\section{Experimental Setup}

\subsection{Use Cases}

We evaluate four explanatory use cases.

\paragraph{UC1. Failure diagnosis and rehydration-point recovery.}
Given a graph that implemented a task incorrectly, can the system isolate the
suspect relationships and identify the upstream node from which context should
be rehydrated for a corrected retry?

\paragraph{UC2. Why-was-this-implemented-like-this analysis.}
Given a task that was implemented in a specific way, can the system
reconstruct the reason for that implementation choice from graph context
alone?

\paragraph{UC3. Interrupted handoff and resumable execution.}
Given a task interrupted after a blocked execution step, can the system
recover why the task was handed off and from which upstream node execution
should resume?

\paragraph{UC4. Constraint-preserving retrieval under token pressure.}
Given a task chosen under a binding safety constraint, can the system preserve
the dominant reason when the rendered context is restricted to a small token
budget?

\subsection{Variants}

UC1--UC3 are executed under three variants: full explanatory context with
detail, full explanatory context without detail, and structural-only context
with detail. UC4 adds a budgeted comparison: a full explanatory baseline at
\mbox{4096} tokens, a full explanatory variant at \mbox{192} tokens, a full
explanatory variant at \mbox{96} tokens, and a structural-only variant at
\mbox{96} tokens. This isolates the contribution of explanatory relationships
from the contribution of extended node detail and tests whether those
relationships survive bounded retrieval.

\subsection{Metrics}

The current harness records explanation roundtrip fidelity, detail roundtrip
fidelity, causal reconstruction score, rendered token count,
rehydration-point hit, retry-success hit, retry-success rate,
dominant-reason hit, and suspect relationship count.

\subsection{Reproducible Artifact}

The paper artifact is generated by running \code{scripts/ci/integration-paper-use-cases.sh}
with \code{bash}. It writes \code{summary.json}, \code{results.md}, \code{results.csv}, and
\code{results-figures.md} under \code{artifacts/paper-use-cases/}.

\section{Results}

Table~\ref{tab:main-results} summarizes the current explanatory-use-case
results. The primary table keeps only the main comparison used in the paper's
central claim. Detail-only and meso-scale ablations are moved to
Appendix~\ref{sec:extended-artifact-results}.

\begin{table*}[t]
\centering
\small
\begin{tabular}{llrrrrrrr}
\toprule
Use case & Variant & Budget & Expl. fid. & Detail fid. & Causal & Retry & Aux. hit & Tokens \\
\midrule
UC1 & full+detail & 4096 & 1.000 & 1.000 & 1.000 & 1 & 1 & 282 \\
UC1 & full-no-detail & 4096 & 1.000 & 0.000 & 0.857 & 1 & 1 & 252 \\
UC1 & detail-only & 4096 & 0.000 & 1.000 & 0.429 & 0 & 0 & 245 \\
UC1 & structural & 4096 & 0.000 & 1.000 & 0.143 & 0 & 0 & 200 \\
UC2 & full+detail & 4096 & 1.000 & 1.000 & 1.000 & n/a & n/a & 285 \\
UC2 & full-no-detail & 4096 & 1.000 & 0.000 & 0.857 & n/a & n/a & 255 \\
UC2 & detail-only & 4096 & 0.000 & 1.000 & 0.429 & n/a & n/a & 247 \\
UC2 & structural & 4096 & 0.000 & 1.000 & 0.143 & n/a & n/a & 203 \\
UC3 & full+detail & 4096 & 1.000 & 1.000 & 1.000 & 1 & 1 & 575 \\
UC3 & full-no-detail & 4096 & 1.000 & 0.000 & 0.857 & 1 & 1 & 535 \\
UC3 & detail-only & 4096 & 0.000 & 1.000 & 0.429 & 0 & 0 & 441 \\
UC3 & structural & 4096 & 0.000 & 1.000 & 0.143 & 0 & 0 & 374 \\
UC4 & full+detail & 4096 & 1.000 & 1.000 & 1.000 & n/a & 1 & 223 \\
UC4 & full@192 & 192 & 1.000 & 0.000 & 1.000 & n/a & 1 & 175 \\
UC4 & full@96 & 96 & 1.000 & 0.000 & 1.000 & n/a & 1 & 89 \\
UC4 & structural@96 & 96 & 0.000 & 0.000 & 0.125 & n/a & 0 & 73 \\
\bottomrule
\end{tabular}
\caption{Explanatory-use-case results from the current repository artifact.
`Retry' reports corrected retry or resume success for UC1 and UC3. `Aux. hit'
reports rehydration or continuation-point hit for UC1 and UC3, and
dominant-reason hit for UC4.}
\label{tab:main-results}
\end{table*}

\subsection{Explanatory Relations Carry the Dominant Signal}

The cleanest result is the size and consistency of the degradation. In UC1,
UC2, and UC3, removing explanation while keeping detail reduces causal
reconstruction from \mbox{1.0} to \mbox{0.143}. Removing detail while
preserving explanation reduces it only to \mbox{0.857}. The dominant loss is
therefore not absence of detail. It is absence of relationship explanation.

\subsection{Detail Improves Completeness, Not Primary Intent}

Detail still matters. The repeated drop from \mbox{1.0} to \mbox{0.857} in
UC1--UC3 shows that explanatory relations do not fully replace node detail.
However, detail behaves as a completeness amplifier rather than the primary
carrier of intent. The same pattern appears in UC4: at a \mbox{192}-token
budget, the explanatory variant loses detail fidelity entirely while
preserving both causal reconstruction and the dominant reason. This suggests
that token budgets should prefer preserving explanatory edges before
preserving every extended detail block.

\subsection{Operational Continuation Depends on Explanation Fields}

UC1 and UC3 show the operational value of the model. In both cases the
explanatory variants recover the correct upstream continuation point, while
the structural-only variants do not. That is strong evidence that fields such
as \texttt{decision\_id} and \texttt{caused\_by\_node\_id} are carrying the
restart or handoff signal.

The closed-loop signal is sharper than continuation hit alone. In both UC1
and UC3, the explanatory variants reach \mbox{1.0} retry-success rate,
including the no-detail setting. Structural-only and detail-only variants stay
at \mbox{0.0}. This matters because detail-only context can still recover part
of the rationale text, but it does not expose a stable machine-readable anchor
from which the system can restart or resume execution.

\subsection{Bounded Retrieval Preserves Reasons When Explanation Survives}

UC4 provides the strongest bounded-retrieval result. The full explanatory
context rendered at \mbox{192} tokens still reaches \mbox{1.0} causal
reconstruction and preserves the dominant reason, even though detail fidelity
falls to \mbox{0.0}. The structural-only variant at \mbox{96} tokens drops to
\mbox{0.125}. This is exactly the systems behavior we want: when budget forces
truncation, explanation should survive before detail.

\subsection{Shorter Context Can Be Worse Context}

Structural-only context is shorter, but materially less useful. In UC1 the
rendered context falls from 282 to 200 tokens. In UC2 it falls from 285 to 203.
In UC3 it falls from 575 to 374. In UC4 the explanatory \mbox{@192} variant
uses 175 tokens, while the structural-only \mbox{@96} variant uses 73. In all
cases that shorter context coincides with severe degradation of causal
reconstruction or loss of the dominant reason.

\subsection{LLM-in-the-Loop Evaluation with Frontier Models}

To validate that the kernel's explanatory context carries signal useful to
actual language models, we run an 18-configuration benchmark with GPT-5.4
(OpenAI) as the inference model and Claude Opus~4 (Anthropic) as an
LLM-as-judge. The benchmark matrix crosses three graph scales (micro: 4~nodes,
meso: 21~nodes, stress: 49~nodes), two domains (operations incident response,
software debugging), and three relation mixes (explanatory, structural, mixed).
For each configuration the kernel generates a synthetic graph, publishes it as
NATS projection events, rehydrates the full bundle, and submits the rendered
context to GPT-5.4 with a structured prompt. Claude Opus~4 then evaluates the
response against domain-aware ground truth on three criteria: task
identification, restart-point selection, and rationale preservation.

\begin{table}[h]
\centering\small
\begin{tabular}{llccc}
\toprule
Scale & Domain & Expl. & Struct. & Mixed \\
\midrule
Micro & Ops    & 3/3 & 2/3 & 3/3 \\
Micro & Debug  & 3/3 & 2/3 & 3/3 \\
Meso  & Ops    & 2/3 & 2/3 & 3/3 \\
Meso  & Debug  & 3/3 & 2/3 & 3/3 \\
Stress & Ops   & 3/3 & 2/3 & 3/3 \\
Stress & Debug & 3/3 & 1/3 & 3/3 \\
\bottomrule
\end{tabular}
\caption{Frontier model benchmark (GPT-5.4 + Claude Opus~4 judge). Each cell
reports TaskOK + RestartOK + ReasonPreserved out of~3.}
\label{tab:frontier}
\end{table}

Explanatory bundles score 17/18 (94\%), structural 11/18 (61\%). Structural
variants consistently lose \texttt{ReasonPreserved} because there is no
rationale metadata to preserve. The judge is calibrated to distinguish
preservation (citing context) from inference (inventing plausible reasons).

To test robustness, we run the same benchmark with a local Qwen3-8B model
(vLLM) and with Claude Opus~4 as inference in cross-validation:

\begin{table}[h]
\centering\small
\begin{tabular}{lccc}
\toprule
Model & Expl. & Struct. & Gap \\
\midrule
GPT-5.4          & 17/18 (94\%) & 11/18 (61\%) & 33pp \\
Claude Opus~4    & 15/18 (83\%) & 10/18 (56\%) & 27pp \\
Qwen3-8B (local) & 18/18 (100\%) & 9/18 (50\%) & 50pp \\
\bottomrule
\end{tabular}
\caption{Cross-model comparison. With explanatory context all models converge
to 83--100\%. Without it, structural accuracy degrades further with smaller
models (61\% $\to$ 50\%). The gap widens from 27pp to 50pp.}
\label{tab:crossmodel}
\end{table}

With explanatory context, all three models converge to near-perfect scores.
Without it, the structural-only gap widens as model capability decreases.

Under token pressure (512-token budget, forcing 68--85\% truncation on
meso and stress graphs), Qwen3-8B drops from 100\% to 83\% on explanatory
bundles --- still matching Claude Opus~4 at full budget. At full budget the
local model achieves 100\% because rationale text is present verbatim and the
model extracts rather than reasons. Under truncation, the kernel's salience
ordering (causal $>$ motivational $>$ evidential $>$ structural) determines
which relationships survive, and the 83\% result reflects genuine context
quality rather than trivial extraction.

\section{Discussion}

The main implication of the current artifact is narrower and more useful than a
generic claim about graph retrieval. The evidence supports a systems claim
about preservation and reuse of explanatory context: when transition-level
explanations are retained across projection, storage, query, and rendering
boundaries, the resulting context is substantially more useful for diagnosis
and recovery than context built from structure alone.

This is also the right DDD split. The kernel does not decide why an action was
taken and does not infer motivation on its own. That meaning is emitted by the
agent or domain application. The kernel's responsibility is to preserve that
meaning in a form that can be rehydrated later and consumed by another runtime.

The result matters because many agent systems currently treat context assembly
as an implementation detail of a single runtime. Our artifact suggests that it
is better treated as a reusable service boundary. Once that boundary exists,
diagnostic and recovery workflows can be evaluated directly instead of being
hidden inside prompt construction logic.

\section{Threats to Validity}

The current evidence has several limitations. First, while the
LLM-in-the-loop benchmark covers micro (4~nodes), meso (21~nodes), and stress
(49~nodes) scales across two domains, these remain synthetic graphs designed to
isolate the explanatory-variable effect. Real-world operational graphs are
likely denser and noisier. Second, the frontier model evaluation uses a single
inference model (GPT-5.4) and a single judge (Claude Opus~4); cross-model
validation would strengthen the result. Third, the current metrics center on
explanatory use cases rather than the full pull-driven and event-driven suite.
Fourth, runtime diversity remains limited. Fifth, explanation quality is
application-supplied; the kernel preserves typed shape and transport, but it
does not validate truthfulness or completeness of emitted rationale. This means
the present results are best interpreted as evidence about preservation and
utility of explanatory context, not about automatic discovery of causal
structure.

Taken together, these limitations bound the claim. The current paper is best
positioned as artifact-backed systems evidence rather than as a large-scale
task benchmark.

\section{Conclusion}

\texttt{rehydration-kernel} supports a concrete systems claim: bounded context
rehydration can be isolated into a reusable kernel, and typed explanatory
relationships materially improve the usefulness of that context for diagnosis,
justification, and recovery.

Across four explanatory use cases, the dominant signal is carried by
relationship explanation rather than by node detail alone. When explanation is
removed, causal reconstruction collapses; when detail is removed but
explanation is preserved, most of the useful trace remains. In operational
cases, explanatory context recovers the correct rehydration or continuation
point. Under token pressure, explanatory context preserves the dominant reason
even after detail is dropped.

The frontier model benchmark confirms this at scale: GPT-5.4 with explanatory
bundles achieves 94\% on task, restart, and rationale metrics across 18
configurations, while structural-only bundles reach 61\%. The gap widens with
graph size, suggesting that the kernel's explanatory context substitutes for
inference-time reasoning that would otherwise need to be performed by the
model. The next step is cross-model validation, noise injection experiments,
and closed-loop evaluation where the selected rehydration point leads to a
corrected retry.

\bibliography{references}

\appendix
\section{Extended Artifact Results}
\label{sec:extended-artifact-results}

\begin{figure*}[t]
\centering
\small
\begin{minipage}{0.47\textwidth}
\centering
\textbf{Detail-only Ablation}\\[0.4ex]
\begin{tabular}{lccc}
\toprule
Case & Causal & Retry & Tokens \\
\midrule
UC1 & 0.429 \metricbar{4.3} & 0 & 245 \\
UC2 & 0.429 \metricbar{4.3} & n/a & 247 \\
UC3 & 0.429 \metricbar{4.3} & 0 & 441 \\
\bottomrule
\end{tabular}
\end{minipage}\hfill
\begin{minipage}{0.47\textwidth}
\centering
\textbf{Meso Robustness}\\[0.4ex]
\begin{tabular}{lccc}
\toprule
Variant & Causal & Retry & Tokens \\
\midrule
UC1 micro & 1.000 \metricbar{10.0} & 1 & 282 \\
UC1 meso & 1.000 \metricbar{10.0} & 1 & 282 \\
\bottomrule
\end{tabular}
\end{minipage}
\caption{Extended artifact ablations omitted from the main table for
readability. `detail-only' preserves partial why-trace in node detail but does
not recover a usable retry anchor. `meso' keeps the same operational result for
UC1 under a denser noisy graph.}
\label{tab:extended-results}
\end{figure*}

\end{document}
